\extrachap{Introduction to Part I}
\label{intro:potential}
\addcontentsline{toc}{chapter}{Introduction to Part I}

The Laplacian describes a potential through its distributional derivatives, mean values and singularity measure. These descriptions support different operations: weak limits, pointwise comparison and analysis of poles. This part proves the passages between them and develops the control of exceptional sets needed for limits and envelopes.

Measures and distributions come first because later curvature currents have measure coefficients. Positive-functional representation identifies those coefficients, convolution supplies smooth approximations, and compactness selects weak limits. Harmonic functions provide the comparison objects. The Poisson kernel solves the Dirichlet problem on a ball; Perron's construction and barriers separate interior existence from attainment of boundary data.

Subharmonic functions are obtained by replacing equality of mean values with an inequality. Their distributional Laplacians are positive measures, and their pointwise values are determined by their almost-everywhere classes. This correspondence is indispensable for upper envelopes: taking a supremum and taking its upper semicontinuous regularization need not give the same values. Riesz decomposition separates a harmonic part from a kernel potential, linking comparison to energy.

Capacity measures the size of a set in terms of the potentials it can support. A set of zero Lebesgue measure can have positive capacity, whereas a polar set cannot carry a nonzero finite-energy measure. Equilibrium potentials and reduction connect this distinction to comparison. The fine topology makes subharmonic functions continuous, and balayage moves a measure while preserving the prescribed potential on the swept set. These constructions explain why exceptional sets are removed only up to capacity zero, rather than by an arbitrary almost-everywhere convention.

The final chapter addresses the transition to several complex variables. Bounds along individual complex lines must be made uniform before linewise subharmonicity can yield a pointwise plurisubharmonic function. An energy estimate, Fourier transform and projection argument provide that uniformity. The same real-variable methods also supply the weak Hessian estimates used later for nonsmooth obstacles. These two estimates have distinct uses in Part II: the incidence bound fixes the pointwise representative, and Hessian regularity controls the determinant of an envelope below a nonsmooth obstacle.
